% ============================================================================
% SUPPLEMENTARY MATERIALS: POLYPHARMACY RISK ASSESSMENT AND ALTERNATIVE THERAPY
% ============================================================================

\documentclass[11pt,a4paper]{article}
\usepackage[utf8]{inputenc}
\usepackage{amsmath,amssymb}
\usepackage{booktabs}
\usepackage{geometry}
\usepackage{graphicx}
\usepackage{hyperref}
\usepackage{float}
\usepackage{longtable}
\usepackage{xcolor}
\geometry{margin=1in}

\hypersetup{
    colorlinks=true,
    linkcolor=blue,
    filecolor=magenta,
    urlcolor=cyan,
}

\title{\textbf{Supplementary Materials}\\
\large Polypharmacy Risk Assessment and Alternative Therapy Suggestion}
\author{Drug-Drug Interaction Knowledge Graph Project}
\date{}

\begin{document}

\maketitle

\tableofcontents
\newpage

% ============================================================================
\section{Polypharmacy Risk Index (PRI) Derivation}
% ============================================================================

\subsection{Network-Based Risk Metrics}

The Polypharmacy Risk Index integrates four complementary network centrality measures that capture different aspects of drug interaction risk within the DDI knowledge graph.

\subsubsection{Degree Centrality}

Degree centrality captures the raw interaction burden of a drug:

\begin{equation}
C_{\text{degree}}(d) = \frac{\text{deg}(d)}{N - 1}
\end{equation}

where $\text{deg}(d)$ is the number of edges connected to drug $d$ in the DDI network and $N$ is the total number of drugs. This metric identifies drugs with high interaction potential.

\subsubsection{Weighted Degree Centrality}

Weighted degree incorporates severity information:

\begin{equation}
C_{\text{weighted}}(d) = \frac{\sum_{j \in \Gamma(d)} w_{dj}}{\sum_{i,j} w_{ij}}
\end{equation}

where $\Gamma(d)$ is the set of neighbors of drug $d$ and $w_{dj}$ is the severity weight of the interaction between $d$ and $j$ (Contraindicated: 4.0, Major: 3.0, Moderate: 2.0, Minor: 1.0).

\subsubsection{Betweenness Centrality}

Betweenness centrality identifies drugs that serve as bridges in risk propagation:

\begin{equation}
C_{\text{betweenness}}(d) = \sum_{s \neq d \neq t} \frac{\sigma_{st}(d)}{\sigma_{st}}
\end{equation}

where $\sigma_{st}$ is the number of shortest paths between drugs $s$ and $t$, and $\sigma_{st}(d)$ is the number of those paths passing through drug $d$.

\subsubsection{Severity Profile Score}

The severity profile quantifies the proportion of high-risk interactions:

\begin{equation}
S(d) = \frac{|\{j : (d,j) \in E \land \text{severity}(d,j) \in \{\text{Contraindicated}, \text{Major}\}\}|}{|\{j : (d,j) \in E\}|}
\end{equation}

\subsection{Weight Optimization}

The PRI weights (0.25, 0.30, 0.20, 0.25) were optimized through grid search to maximize correlation with clinical severity assessments while maintaining balanced contribution from each metric. The weighted degree component receives the highest weight (0.30) as it most directly captures clinically significant interactions.

\subsection{PRI Distribution Statistics}

\begin{table}[H]
\centering
\caption{PRI Score Distribution Statistics Across All Drugs}
\label{tab:pri_stats}
\begin{tabular}{lr}
\toprule
\textbf{Statistic} & \textbf{Value} \\
\midrule
Mean & 0.342 \\
Median & 0.298 \\
Standard Deviation & 0.186 \\
Interquartile Range & 0.245 \\
Skewness & 0.824 (right-skewed) \\
Kurtosis & 2.91 \\
Minimum & 0.000 \\
Maximum & 0.892 \\
\bottomrule
\end{tabular}
\end{table}

% ============================================================================
\section{High-Risk Drug Categories}
% ============================================================================

\subsection{Top 10 Drugs by Interaction Count}

\begin{table}[H]
\centering
\caption{Drugs with Highest Interaction Counts in the Knowledge Graph}
\label{tab:high_pri_drugs}
\begin{tabular}{llccr}
\toprule
\textbf{Drug} & \textbf{Category} & \textbf{Total DDIs} & \textbf{Severe} & \textbf{Severe \%} \\
\midrule
Quinidine & Antiarrhythmic & 4,782 & 258 & 5.4\% \\
Procaine & Local Anesthetic & 4,719 & 786 & 16.7\% \\
Lidocaine & Local Anesthetic & 4,208 & 626 & 14.9\% \\
Verapamil & Calcium Channel Blocker & 4,208 & 384 & 9.1\% \\
Propranolol & Beta-blocker & 4,184 & 208 & 5.0\% \\
Warfarin & Anticoagulant & 4,145 & 298 & 7.2\% \\
Isradipine & Calcium Channel Blocker & 4,142 & 162 & 3.9\% \\
Pindolol & Beta-blocker & 4,048 & 722 & 17.8\% \\
Indomethacin & NSAID & 4,045 & 778 & 19.2\% \\
Nicardipine & Calcium Channel Blocker & 4,044 & 402 & 9.9\% \\
\bottomrule
\end{tabular}
\end{table}

\subsection{Risk by Drug Category}

\begin{table}[H]
\centering
\caption{Average PRI Scores by Drug Category}
\label{tab:category_pri}
\begin{tabular}{lcc}
\toprule
\textbf{Drug Category} & \textbf{Mean PRI} & \textbf{Drug Count} \\
\midrule
Anticoagulants & 0.723 & 28 \\
Antiarrhythmics & 0.698 & 42 \\
Anticonvulsants & 0.654 & 67 \\
Immunosuppressants & 0.612 & 35 \\
Antifungals (Azoles) & 0.589 & 23 \\
HIV Antivirals & 0.567 & 48 \\
Antipsychotics & 0.445 & 89 \\
Antidepressants & 0.398 & 134 \\
Antihypertensives & 0.312 & 156 \\
NSAIDs & 0.287 & 78 \\
\bottomrule
\end{tabular}
\end{table}

% ============================================================================
\section{Multi-Objective Recommendation Framework}
% ============================================================================

\subsection{Therapeutic Similarity Calculation}

The therapeutic similarity score combines multiple dimensions of drug equivalence:

\begin{equation}
T(d_{\text{alt}}) = \alpha_1 \cdot \text{ATC}(d_{\text{alt}}) + \alpha_2 \cdot \text{Target}(d_{\text{alt}}) + \alpha_3 \cdot \text{Disease}(d_{\text{alt}}) + \alpha_4 \cdot \text{Pathway}(d_{\text{alt}})
\end{equation}

where $\alpha_1 = 0.35$, $\alpha_2 = 0.25$, $\alpha_3 = 0.25$, $\alpha_4 = 0.15$.

\subsubsection{ATC Code Matching Levels}

\begin{table}[H]
\centering
\caption{ATC Code Matching Score by Level}
\label{tab:atc_levels}
\begin{tabular}{llc}
\toprule
\textbf{ATC Level} & \textbf{Description} & \textbf{Score} \\
\midrule
5th Level & Same chemical substance & 1.00 \\
4th Level & Same chemical subgroup & 0.80 \\
3rd Level & Same pharmacological subgroup & 0.50 \\
2nd Level & Same therapeutic subgroup & 0.20 \\
1st Level & Same anatomical main group & 0.05 \\
None & No match & 0.00 \\
\bottomrule
\end{tabular}
\end{table}

\subsubsection{Protein Target Overlap}

Target similarity is computed using Jaccard index:

\begin{equation}
\text{Target}(d_{\text{alt}}) = \frac{|T_{d_{\text{orig}}} \cap T_{d_{\text{alt}}}|}{|T_{d_{\text{orig}}} \cup T_{d_{\text{alt}}}|}
\end{equation}

where $T_d$ represents the set of protein targets for drug $d$.

\subsubsection{Disease Indication Coverage}

Disease similarity measures how well the alternative covers the original drug's indications:

\begin{equation}
\text{Disease}(d_{\text{alt}}) = \frac{|D_{d_{\text{orig}}} \cap D_{d_{\text{alt}}}|}{|D_{d_{\text{orig}}}|}
\end{equation}

\subsection{Safety Score Computation}

\begin{table}[H]
\centering
\caption{Safety Penalty Weights by Interaction Severity}
\label{tab:safety_penalties}
\begin{tabular}{lcc}
\toprule
\textbf{Severity} & \textbf{Penalty Weight} & \textbf{Normalized} \\
\midrule
Contraindicated & 4.0 & 1.000 \\
Major & 3.0 & 0.750 \\
Moderate & 2.0 & 0.500 \\
Minor & 1.0 & 0.250 \\
\bottomrule
\end{tabular}
\end{table}

% ============================================================================
\section{Case Study Extended Analysis}
% ============================================================================

\subsection{Complete DDI List for Original Regimen}

\begin{longtable}{llp{7cm}}
\caption{Complete DDI Analysis for High-Risk Cardiovascular Regimen (from Knowledge Graph)} \label{tab:ddi_details} \\
\toprule
\textbf{Drug Pair} & \textbf{Severity} & \textbf{Mechanism} \\
\midrule
\endfirsthead
\multicolumn{3}{c}{\tablename\ \thetable{} -- continued} \\
\toprule
\textbf{Drug Pair} & \textbf{Severity} & \textbf{Mechanism} \\
\midrule
\endhead
\bottomrule
\endfoot
Warfarin + Amiodarone & Major & Amiodarone may increase risk or severity of bleeding when combined with Warfarin \\
Warfarin + Digoxin & Moderate & Digoxin may decrease excretion rate of Warfarin \\
Warfarin + Quinidine & Moderate & Quinidine may increase serum concentration of Warfarin \\
Warfarin + Propranolol & Moderate & Propranolol may increase serum concentration of Warfarin \\
Amiodarone + Digoxin & Moderate & Amiodarone may increase serum concentration of Digoxin \\
Amiodarone + Quinidine & Moderate & Amiodarone may increase serum concentration of Quinidine \\
Amiodarone + Propranolol & Moderate & Amiodarone may increase therapeutic efficacy of Propranolol \\
Digoxin + Quinidine & Moderate & Quinidine may increase serum concentration of Digoxin \\
Digoxin + Propranolol & Moderate & Propranolol may increase arrhythmogenic activity of Digoxin \\
Quinidine + Propranolol & Moderate & Quinidine may increase serum concentration of Propranolol \\
\end{longtable}

\subsection{Detailed Alternative Scoring Matrix}

\begin{table}[H]
\centering
\caption{Detailed Scoring Matrix for Warfarin Alternatives}
\label{tab:detailed_scoring}
\begin{tabular}{@{}lccccccc@{}}
\toprule
\textbf{Alternative} & \textbf{ATC} & \textbf{Target} & \textbf{Disease} & \textbf{Pathway} & \textbf{$T$} & \textbf{SAF} & \textbf{$R$} \\
\midrule
Dabigatran & 0.80 & 0.92 & 0.95 & 0.72 & 0.85 & 0.78 & 0.75 \\
Fondaparinux & 0.80 & 0.88 & 0.90 & 0.78 & 0.85 & 0.35 & 0.25 \\
Apixaban & 0.80 & 0.90 & 0.92 & 0.75 & 0.85 & 0.13 & 0.00 \\
Rivaroxaban & 0.80 & 0.89 & 0.91 & 0.74 & 0.85 & 0.13 & 0.00 \\
Edoxaban & 0.80 & 0.88 & 0.90 & 0.73 & 0.85 & 0.13 & 0.00 \\
\bottomrule
\end{tabular}
\end{table}

\subsection{Post-Substitution DDI Analysis}

\begin{table}[H]
\centering
\caption{DDIs After Warfarin $\rightarrow$ Dabigatran Substitution}
\label{tab:post_sub_ddi}
\begin{tabular}{lll}
\toprule
\textbf{Drug Pair} & \textbf{Severity} & \textbf{Status} \\
\midrule
\multicolumn{3}{l}{\textit{Unchanged interactions (6):}} \\
Amiodarone + Digoxin & Moderate & Unchanged \\
Amiodarone + Quinidine & Moderate & Unchanged \\
Amiodarone + Propranolol & Moderate & Unchanged \\
Digoxin + Quinidine & Moderate & Unchanged \\
Digoxin + Propranolol & Moderate & Unchanged \\
Quinidine + Propranolol & Moderate & Unchanged \\
\midrule
\multicolumn{3}{l}{\textit{New interactions with Dabigatran (4):}} \\
Dabigatran + Amiodarone & Moderate & New (replaces Major) \\
Dabigatran + Digoxin & Moderate & New \\
Dabigatran + Quinidine & Moderate & New \\
Dabigatran + Propranolol & Moderate & New \\
\midrule
\multicolumn{3}{l}{\textit{Eliminated interactions (4):}} \\
Warfarin + Amiodarone & -- & Eliminated (was Major) \\
Warfarin + Digoxin & -- & Eliminated (was Moderate) \\
Warfarin + Quinidine & -- & Eliminated (was Moderate) \\
Warfarin + Propranolol & -- & Eliminated (was Moderate) \\
\bottomrule
\end{tabular}
\end{table}

% ============================================================================
\section{Algorithm Pseudocode}
% ============================================================================

\subsection{PRI Computation Algorithm}

\begin{verbatim}
Algorithm: Polypharmacy Risk Index Computation

Input: 
  - G: DDI Knowledge Graph
  - drug_regimen: Set of drugs to evaluate

Output:
  - pri_scores: Dictionary mapping drug -> PRI score

1. for each drug d in drug_regimen:
2.     deg = count_neighbors(d, G)
3.     weighted_deg = sum(severity_weight(d, j) for j in neighbors(d))
4.     betweenness = compute_betweenness(d, G)
5.     severe_count = count_severe_interactions(d, G)
6.     total_interactions = count_all_interactions(d, G)
7.     
8.     C_degree = normalize(deg)
9.     C_weighted = normalize(weighted_deg)
10.    C_betweenness = normalize(betweenness)
11.    S = severe_count / total_interactions if total_interactions > 0 else 0
12.    
13.    pri_scores[d] = 0.25*C_degree + 0.30*C_weighted + 
                       0.20*C_betweenness + 0.25*S
14. 
15. return pri_scores
\end{verbatim}

\subsection{Multi-Objective Recommendation Algorithm}

\begin{verbatim}
Algorithm: Multi-Objective Drug Recommendation

Input: 
  - current_regimen: Set of drugs
  - target_drug: Drug to be replaced
  - kg: Knowledge graph

Output:
  - ranked_alternatives: List of (drug, score) tuples

1. candidates = find_therapeutic_alternatives(target_drug, kg)
2. for each candidate in candidates:
3.     # Compute therapeutic similarity
4.     atc_score = compute_atc_similarity(candidate, target_drug)
5.     target_score = compute_target_overlap(candidate, target_drug, kg)
6.     disease_score = compute_disease_overlap(candidate, target_drug, kg)
7.     pathway_score = compute_pathway_similarity(candidate, target_drug, kg)
8.     T = 0.35*atc_score + 0.25*target_score + 
          0.25*disease_score + 0.15*pathway_score
9.     
10.    # Compute safety improvement
11.    original_penalty = sum_severity_penalties(target_drug, current_regimen)
12.    candidate_penalty = sum_severity_penalties(candidate, current_regimen)
13.    SAF = 1 - (candidate_penalty / original_penalty)
14.    
15.    # Compute risk reduction
16.    delta_contra = count_contraindicated(target_drug) - 
                      count_contraindicated(candidate)
17.    delta_major = count_major(target_drug) - count_major(candidate)
18.    delta_pri = compute_pri_change(target_drug, candidate, current_regimen)
19.    R = normalize(delta_contra + delta_major + 0.5*delta_pri)
20.    
21.    # Compute final score
22.    score = 0.40*T + 0.35*SAF + 0.25*R
23.    add (candidate, score) to ranked_alternatives
24. 
25. sort ranked_alternatives by score descending
26. return ranked_alternatives
\end{verbatim}

% ============================================================================
\section{Validation Against Clinical Guidelines}
% ============================================================================

The PRI thresholds and recommendation system were validated against established clinical guidelines:

\subsection{Reference Guidelines}

\begin{itemize}
    \item \textbf{AGS Beers Criteria\textsuperscript{\textregistered} 2019}: Potentially inappropriate medications in older adults
    \item \textbf{STOPP/START Criteria v2}: Screening tool for medication review in older persons
    \item \textbf{NHS Scotland Polypharmacy Guidance 2018}: Evidence-based approach to polypharmacy management
    \item \textbf{HEDIS Medication Reconciliation}: Quality measures for medication management
\end{itemize}

\subsection{Validation Metrics}

\begin{table}[H]
\centering
\caption{Recommendation System Validation Results}
\label{tab:validation}
\begin{tabular}{lc}
\toprule
\textbf{Metric} & \textbf{Value} \\
\midrule
Average Risk Reduction & 28.4\% \\
Therapeutic Equivalence Maintained & 94.2\% \\
Mean Recommendation Score & 0.623 \\
Top-1 Clinical Acceptance Rate & 78.5\% \\
Top-3 Clinical Acceptance Rate & 92.1\% \\
PRI Threshold Alignment with Guidelines & 87.3\% \\
\bottomrule
\end{tabular}
\end{table}

% ============================================================================
\section*{References}
% ============================================================================

\bibliographystyle{plain}
\begin{thebibliography}{6}
\bibitem{wishart2018drugbank} Wishart DS, et al. DrugBank 5.0: a major update to the DrugBank database for 2018. \textit{Nucleic Acids Res}. 2018;46(D1):D1074-D1082.
\bibitem{xiong2022ddinter} Xiong G, et al. DDInter: an online drug-drug interaction database. \textit{Nucleic Acids Res}. 2022;50(D1):D1200-D1207.
\bibitem{beers2019} American Geriatrics Society 2019 Beers Criteria Update Expert Panel. American Geriatrics Society 2019 Updated AGS Beers Criteria. \textit{J Am Geriatr Soc}. 2019;67(4):674-694.
\bibitem{stopp2015} O'Mahony D, et al. STOPP/START criteria for potentially inappropriate prescribing in older people: version 2. \textit{Age Ageing}. 2015;44(2):213-218.
\bibitem{nhs2018} NHS Scotland. Polypharmacy Guidance, Realistic Prescribing. 3rd ed. Edinburgh: Scottish Government; 2018.
\bibitem{newman2010networks} Newman MEJ. Networks: An Introduction. Oxford University Press; 2010.
\end{thebibliography}

\end{document}
